% =====================================================================
%  5.12 TECNOLOGIAS E INTEGRACION DE PRODUCTOS DE TERCEROS
%  Rubrica: describirlas PORMENORIZADAMENTE y justificarlas CON DETALLE.
%  Ojo: hay que describir explicitamente lo que NO es de tu autoria y
%  justificar por que lo has usado.
% =====================================================================

\chapter{Tecnologías y productos de terceros}
\label{cap:tecnologias}

\section{Lenguaje y framework}

% Version concreta y motivo de la eleccion.

\section{Bibliotecas de terceros}

\begin{table}[H]
  \centering
  \begin{tabular}{@{}llp{6cm}@{}}
    \toprule
    \textbf{Producto} & \textbf{Versión} & \textbf{Uso y justificación} \\
    \midrule
     &  &  \\
     &  &  \\
    \bottomrule
  \end{tabular}
  \caption{Productos de terceros integrados en el sistema.}
  \label{tab:terceros}
\end{table}

\section{Herramientas de desarrollo}

% Control de versiones, entorno, linters, CI si la hay.

\section{Fuentes de datos clínicos}

% La guia STOPP/START v3 como fuente. Importante dejar clara la
% procedencia y que no es de autoria propia.

\section{Licencias}

% Licencia de cada producto de terceros y compatibilidad entre ellas.
